% ============================================================
% Causal Memory: Formulation v0.1 -- FOUNDATIONS (share version)
% The ~3-page slice for step-by-step verification with Yujia.
% The full draft (Layer-2 analysis, selection section, E0/E1
% pre-registration, positioning) lives in
% causal_memory_formulation.tex and enters at v0.2+ as these
% foundations receive verdicts.
% Compile with pdfLaTeX. Single self-contained file for Overleaf.
% ============================================================
\documentclass[10pt]{article}
\usepackage[margin=0.9in]{geometry}
\usepackage{amsmath,amssymb,amsthm}
\usepackage{xcolor}
\usepackage[colorlinks=true,linkcolor=blue!60!black,citecolor=blue!60!black,urlcolor=blue!60!black]{hyperref}

\newcommand{\indep}{\perp\!\!\!\perp}
\newcommand{\nindep}{\not\!\perp\!\!\!\perp}
\newcommand{\Do}{\mathrm{do}}
\newcommand{\An}{\mathrm{An}}
\newcommand{\Pa}{\mathrm{Pa}}
\newcommand{\past}{V_{\le t}}
\newcommand{\future}{V_{> t}}
\newcommand{\checkY}[1]{\textcolor{teal}{\footnotesize [check: #1]}}

\makeatletter
\def\thm@space@setup{\thm@preskip=4pt plus 1pt \thm@postskip=4pt plus 1pt}
\renewcommand\section{\@startsection{section}{1}{\z@}%
  {-2.2ex \@plus -.6ex \@minus -.2ex}{1.2ex \@plus .2ex}%
  {\normalfont\large\bfseries}}
\makeatother
\AtBeginDocument{%
  \setlength{\abovedisplayskip}{5pt plus 1pt}%
  \setlength{\belowdisplayskip}{5pt plus 1pt}}

\theoremstyle{plain}
\newtheorem{lemma}{Lemma}
\theoremstyle{definition}
\newtheorem{assumption}{A}
\newtheorem{definitionx}{D}
\newtheorem{questionx}{Q}
\newtheorem{remark}{Remark}

\begin{document}
\begin{center}
{\LARGE\bfseries Causal Memory: Formulation Document (v0.1)}\\[3pt]
{\large Foundations for step-by-step verification: setup, definitions, and the frontier lemma}\\[5pt]
{\small Mian (Nothern) Wu \;\textperiodcentered\; prepared for verification with Yujia Zheng \;\textperiodcentered\; July 2026}
\end{center}
\vspace{2pt}

\begin{abstract}
\noindent Memory should be defined as time-delayed dependency structure within the true data-generating process; a memory system is \emph{trustworthy} exactly when its stored values and its write/read structure correspond, up to allowed indeterminacies, to that process. This note contains only the foundations to be checked first: setup (A1--A9), definitions (D1--D3), one proved lemma, and decision points (Q1--Q7). The Layer-2 identification analysis and the pre-registered experiments are drafted separately and enter at v0.2 once these have verdicts; nothing runs before sign-off. Inline \checkY{like this} flags mark where I most want correction. I am committing to the theory track: the identifiability framing is what distinguishes this direction from the memory-agent literature.
\end{abstract}

\section{Thesis}\label{sec:thesis}

An agent should not memorize every signal at every timestamp; it should store what corresponds to the truthful causal data-generating process, and read the past through that process's dependency structure. Two claims make this precise. \emph{Definitional}: memory is time-delayed dependency among the variables (observed or latent) of a temporal causal model, not a separate external module; an external store is just a variable with identity dynamics and gated write/read edges (Remark~\ref{rem:external}). \emph{Epistemic}, where the contribution lies: without identifiability guarantees, a learned memory representation $\hat z$ need not correspond to the true hidden process $z$; causal discovery and causal representation learning supply those guarantees, making \emph{trustworthy memory} (D\ref{def:trust}) the target; reliability and efficiency follow as byproducts.

\paragraph{Causal discovery versus causal representation learning (my understanding, for checking).}
Causal discovery takes variables whose values are given and recovers the \emph{structure} among them: a data table in, a graph over its variables out. Causal representation learning is deeper: the relevant variables are latent, so their \emph{values} must first be recovered from observations $x = g(z)$ up to a permitted indeterminacy, and only then structure over them. Identifiability is the guarantee that the estimated $\hat z$ stands in a controlled (ideally component-wise invertible) relation to the true $z$. Memory needs the CRL level of guarantee at Layer~2: what the agent stores must equal the true latent value, not merely a predictive embedding of it. \checkY{is this restatement of the CD/CRL distinction accurate?}

\paragraph{Why this graph is not a knowledge graph.}
A knowledge graph asserts relations among entities; its edges encode claims. The graph here is the \emph{generating mechanism}: nodes are variables of the process that produces each observation, edges are its direct causal dependencies, and it supports interventional and counterfactual statements, which relation graphs do not.

\section{Formal setup}\label{sec:setup}

Let $z_t = (z_{t,1},\dots,z_{t,n}) \in \mathbb{R}^n$ evolve over $t = 1,\dots,T$ within episodes drawn i.i.d., with unrolled DAG $G$ over $V = \{z_{t,i}\}$ and assignments
\[
z_{t,i} \;=\; f_i\big(\Pa(z_{t,i}),\, \varepsilon_{t,i}\big),
\qquad \Pa(z_{t,i}) \subseteq \{z_{t-\tau, j} : 0 \le \tau \le L\}.
\]

\begin{assumption}\label{a:scm}
The process is a temporal SCM as above with maximum lag $L$ and acyclic instantaneous ($\tau=0$) dependencies.
\end{assumption}

\begin{assumption}\label{a:gated}
(Gated stationarity.) $G$ and $\{f_i\}$ are time-invariant; time-variation enters only through a regime process $u_t \in \mathcal{U}$ (e.g., $\mathsf{write}_j$, $\mathsf{hold}$, $\mathsf{read}_j$) gating edge activations: each edge $e$ carries $\alpha_e(u_t) \in \{0,1\}$, and $f_i$ may depend on $u_t$. \checkY{Q1: this, or a regime-free nonstationary-transition formulation?}
\end{assumption}

\begin{assumption}\label{a:noise}
Noises are mutually independent across components and time; regime-dependent distributions allowed.
\end{assumption}

\begin{assumption}\label{a:obs}
(Observation models.) \textbf{(O1)} Layer 1: $s_t = z_t$ observed. \textbf{(O2a)} $x_t = g(z_t)$, $g$ unknown, injective; intended instance: an agent's internal state, which carries its memory. \textbf{(O2b)} memory and cue are excluded from what generates $x_t$ during holds; per-step inversion is then impossible and identification becomes a filtering problem (Layer-2 note).
\end{assumption}

\begin{assumption}\label{a:regime}
Regime labels $u_t$ are observed, or estimable by change-point detection; known by construction in synthetic settings.
\end{assumption}

\begin{assumption}\label{a:mf}
Markov and faithfulness hold for $(G, P)$. Caveat: near-deterministic holds do not violate faithfulness of the true model, but they destroy the power of tests that rely on it; population and estimation claims are kept separate.
\end{assumption}

\begin{assumption}\label{a:nosel}
No selection on the data, unless explicitly modeled and tested. False by construction for goal-filtered agent corpora, where subgoals are provably selections on $(s_t, a_t)$ \cite{qiu2024subtask}; certificate (iii) of D\ref{def:trust} exists because of this.
\end{assumption}

\begin{assumption}\label{a:policy}
Data come from a fixed behavior policy $\pi$; recovered structure is relative to $P_\pi$. \checkY{Q7: is policy-relativity acceptable in the definition?}
\end{assumption}

\begin{assumption}\label{a:suff}
(Layer 1 only.) Causal sufficiency over the observed $s_t$; relaxable to FCI-class discovery at the cost of weaker output.
\end{assumption}

\paragraph{The write-hold-read motif (unit of memory).}
Fix a cue $c$, a memory latent $m$, a readout $y$, and distractors $d$ with full-rank stationary dynamics; write at $t_0$, read at $t_1 = t_0 + \Delta$:
\[
\underbrace{m_{t_0} = w(c_{t_0}, \varepsilon)}_{\text{write, } u_{t_0} = \mathsf{write}}
\qquad
\underbrace{m_t = m_{t-1} + \sigma\,\varepsilon_t \ \ (t_0 < t < t_1)}_{\text{hold}}
\qquad
\underbrace{y_{t_1} = r(m_{t_1}, d_{t_1}, \varepsilon)}_{\text{read, } u_{t_1} = \mathsf{read}}
\]
The cue varies across episodes; the load-$k$ variant runs $k$ items, simultaneous or staggered. Two knobs organize everything: hold noise $\sigma \ge 0$ and delay $\Delta$. Per our meeting, $T$ is not a theoretical bound; the estimator's window $W$ is an engineering threshold to sweep.

\section{Memory, masks, and trustworthiness}\label{sec:defs}

\begin{definitionx}[Memory carrier]\label{def:carrier}
The process $(m_s)_{t_0 \le s \le t_1}$ is a \emph{memory carrier} for $(c_{t_0}, y_{t_1})$ if
(i) $c_{t_0} \nindep y_{t_1}$ under $P_\pi$ with $\Delta \ge \Delta_{\min}$, and
(ii) for every $s \in (t_0, t_1]$ and constant $\tilde m$, $c_{t_0} \indep y_{t_1}$ under $P_\pi^{\Do(m_s := \tilde m)}$.
Equivalently, every active path from $c_{t_0}$ to $y_{t_1}$ intersects $\{m_s\}$. Memory is thus mediated long-lag dependence through a persistent variable, an interventional notion, which separates storage from selection-induced dependence.
\end{definitionx}

\begin{remark}\label{rem:external}
An external memory module is the special case with exact identity dynamics ($\sigma = 0$) and gated write/read edges implemented by retrieval; internal states and external stores are the same object in $G$.
\end{remark}

\begin{definitionx}[Causal frontier / memory mask]\label{def:frontier}
With $\past$ the nodes at times $\le t$, $Y \subseteq \future$ the task variables in a horizon window, and $\An^*(Y) = \An(Y) \cup Y$:
\[
F_t(Y) \;=\; \big\{\, v \in \past \;:\; \exists\, w \in \future,\ v \to w \in G,\ w \in \An^*(Y) \,\big\}.
\]
\end{definitionx}

\begin{lemma}[$F_t$ is the unique minimum memory]\label{lem:frontier}
Under A\ref{a:scm} and the Markov property, \emph{(i)} $Y \indep \past \setminus F_t \,\big|\, F_t$. Under faithfulness in addition, \emph{(ii)} every $S \subseteq \past$ with $Y \indep \past \setminus S \mid S$ satisfies $F_t \subseteq S$. Hence $F_t$ is the unique inclusion-minimal past-measurable set screening off the future task variables, and $|F_t|$ lower-bounds any lossless memory of the past.
\end{lemma}

\begin{proof}
\emph{(i)} Let $S \subseteq \past$ and consider a path from $a \in \past$ to $y \in Y$. Edges never point backward in time, so the step at which the path first moves from past to future traverses an edge $v \to w$ forward, arrowhead at $w \in \future$. From $w$ onward the path cannot contain a future collider: a future node's descendants are all future, so none lies in $S$, and any collider there blocks the path. Arriving at $w$ with an arrowhead, the path must leave along an outgoing edge, and inductively its suffix is a directed path $w \to \cdots \to y$ inside $\future$; in particular it never re-enters the past. Hence $w \in \An^*(Y)$ and $v \in F_t$. Now take $S = F_t$ and $a \in \past \setminus F_t$: on any such path, $v$ is a non-collider (the path leaves along the tail of $v \to w$), $v \ne a$, and $v \in S$, so the path is blocked. No active path exists, giving (i) by Markov.

\emph{(ii)} Let $f \in F_t$ with boundary edge $f \to w$, $w \in \An^*(Y)$, and pick a directed path $w \to \cdots \to y$, all inside $\future$. If $f \notin S$, the path $f \to w \to \cdots \to y$ is directed (no colliders) with no non-endpoint node in $S \subseteq \past$; it is active given $S$, and faithfulness gives $f \nindep y \mid S$, contradicting separation. So every separating $S$ contains $F_t$, and $F_t$ itself separates by (i).
\end{proof}

\begin{remark}[Consequences]\label{rem:frontier}
(a) \emph{Planning}: $\mathbb{E}[\phi(Y) \mid \past] = \mathbb{E}[\phi(Y) \mid F_t]$, the ``conditional set for planning'' from our meeting made exact; it also answers ``what to remember'' as a corollary.
(b) \emph{Online maintenance (fixed $Y$)}: $F_{t+1}(Y) \subseteq F_t(Y) \cup V_{t+1}$: add the new slice, retire nodes whose forward-relevant children have passed. For a sliding horizon window the task set changes with $t$ and the recursion needs restating. \checkY{is fixed-$Y$ the right scope, or should the sliding-window recursion be worked out?}
(c) \emph{Task conditioning}: $F_t$ depends on $Y$ and the horizon, matching task-conditioned representations \cite{zheng2026specialist}.
(d) \emph{Scope}: a population statement about the true $G$; estimability in memory regimes is the Layer-1/Layer-2 question. \checkY{Q3: happy with $F_t$ as the formal answer to ``what deserves storage''?}
\end{remark}

\begin{definitionx}[Trustworthy memory]\label{def:trust}
An estimated system ($\hat z$, $\hat G$, write/read maps) is \emph{trustworthy} on a memory-relevant index set $M$ if:
\emph{(i) value correspondence}: $\hat z_{\rho(i)} = h_i(z_i)$ for all $i \in M$, for some permutation $\rho$ and component-wise invertible $h_i$; where this is provably impossible (simultaneously frozen items admit an arbitrary invertible mixing mid-hold), the certificate relaxes to block level (Layer-2 note);
\emph{(ii) structure correspondence}: write/hold/read edges of $\hat G$ match $G$ on $M$ (edge $F_1$, SHD);
\emph{(iii) selection soundness}: on data whose long-lag dependence arises from selection alone, the system reports \emph{no carrier} \cite{zheng2024selection}.
Byproducts: with (i)+(ii), interventions on $\hat m$ implement interventions on $m$ (editing specificity, auditable recall, detectable confabulation); by Lemma~\ref{lem:frontier}, retaining $F_t$ shrinks the budget from $O(nT)$ to $|F_t|$ with no task-relevant loss.
\end{definitionx}

\section{The program from here}\label{sec:program}

Layer 1 runs temporal discovery on observed $s_t$ (PC-stable; PCMCI$^+$ \cite{runge2020}; CD-NOD with $u_t$ as context \cite{huang2020}, the natural match to A\ref{a:gated}), uses $\hat F_t(Y)$ as retention mask and planning set, and treats its memory-specific failure modes (near-deterministic holds, selection, policy confounding, latent confounders) as experimental arms. Layer 2 asks when the stored \emph{values} are identifiable: write/read events supply auxiliary variation \cite{yao2021leap}, IDOL-type sparsity covers the ambient process \cite{li2025idol}, frozen blocks are set-identifiable at best mid-hold \cite{zheng2025concepts}, and distinct read children, staggered events, or non-Gaussian hold noise \cite{zheng2022} restore component-wise recovery; precise statements enter at v0.2. A selection gate \cite{zheng2024selection} runs at both layers before any memory edge is asserted. E0 (discovery on the observed motif) and E1 (CRL under unknown mixing) are pre-registered in the full draft, gated on sign-off.

\section{Questions}\label{sec:questions}

\begin{questionx}
Gated SCM (A\ref{a:gated}, A\ref{a:regime}), or a regime-free nonstationary-transition formulation with events implicit?
\end{questionx}
\begin{questionx}
Layer 2: is regime-change boundary identification + set-identifiability within holds + gluing along the diagonal hold transition viable, or is the nonparametric-concepts route cleaner?
\end{questionx}
\begin{questionx}
Is $F_t$ (D\ref{def:frontier}, Lemma~\ref{lem:frontier}) the right formalization of ``memorize only what matters,'' or do you prefer a Markov-blanket or soft information-theoretic variant?
\end{questionx}
\begin{questionx}
Selection: preprocessing gate, or jointly modeled inside the estimator?
\end{questionx}
\begin{questionx}
Layer 1: publishable on its own, or scaffolding for the Layer-2 paper?
\end{questionx}
\begin{questionx}
May small synthetic sanity checks run while this document iterates, or does all compute wait for sign-off?
\end{questionx}
\begin{questionx}
Is $P_\pi$-relative structure (A\ref{a:policy}) acceptable, or should cross-policy invariance enter the definition of memory itself?
\end{questionx}

\begin{thebibliography}{9}\footnotesize\setlength{\itemsep}{1pt}\setlength{\parskip}{0pt}

\bibitem{yao2021leap} W. Yao, Y. Sun, A. Ho, C. Sun, K. Zhang. Learning temporally causal latent processes from general temporal data. ICLR 2022. arXiv:2110.05428.

\bibitem{li2025idol} Z. Li, Y. Shen, K. Zheng, R. Cai, X. Song, M. Gong, G. Chen, K. Zhang. On the identification of temporally causal representation with instantaneous dependence. ICLR 2025. arXiv:2405.15325.

\bibitem{zheng2022} Y. Zheng, I. Ng, K. Zhang. On the identifiability of nonlinear ICA: sparsity and beyond. NeurIPS 2022. arXiv:2206.07751.

\bibitem{zheng2025concepts} Y. Zheng, S. Xie, K. Zhang. Nonparametric identification of latent concepts. ICML 2025.

\bibitem{zheng2024selection} Y. Zheng, Z. Tang, Y. Qiu, B. Sch\"olkopf, K. Zhang. Detecting and identifying selection structure in sequential data. ICML 2024. arXiv:2407.00529.

\bibitem{qiu2024subtask} Y. Qiu, Y. Zheng, K. Zhang. Identifying selections for unsupervised subtask discovery. NeurIPS 2024. arXiv:2410.21616.

\bibitem{zheng2026specialist} Y. Zheng, F. Feng, Y. Li, S. Xie, K. Murphy, K. Zhang. From generalist to specialist representation. ICML 2026. arXiv:2605.12733.

\bibitem{huang2020} B. Huang, K. Zhang, J. Zhang, J. Ramsey, R. Sanchez-Romero, C. Glymour, B. Sch\"olkopf. Causal discovery from heterogeneous/nonstationary data. JMLR 2020.

\bibitem{runge2020} J. Runge. Discovering contemporaneous and lagged causal relations in autocorrelated nonlinear time series datasets. UAI 2020.

\end{thebibliography}

\end{document}
